\section{Committing-Notion Bounds}
\label{app:committing-transfer}

The headline committing results bound full-input CMTDs-4
(\Cref{cor:cmtds4}) and its tag-projected strengthening
(\Cref{thm:tag-cmtds4}). The remaining \cite{BH22} committing
notions---CMTDs-$\ell$ and the CMTs-$\ell$ source games for
$\ell \in \{1,3,4\}$---all reduce to the same root-collision bound, so we
record the two reduction patterns once rather than accounting for each
notion separately.

\begin{lemma}[Forwarding Reductions for the D-Notions]
\label{lem:cmtds-identity-reductions}
Fix $s \geq 2$ and $X \in \{\mathsf{CMTDs\text{-}1},
\mathsf{CMTDs\text{-}4}\}$. The identity reduction
$\mathcal{A}^\ast = \mathcal{A}$---preceded, for
$X = \mathsf{CMTDs\text{-}1}$, by the source-game syntax, length, and
pairwise-$\WiC_1$ prechecks---turns every $s$-way $X$-adversary against
\TW{} into an $s$-way CMTDs-3 adversary with
\[
  \Adv^X_\TW(\mathcal{A}) \leq \Advcmtds{3}_\TW(\mathcal{A}^\ast),
  \quad
  \Nprim(\mathcal{A}^\ast) = \Nprim(\mathcal{A}),
  \quad
  \Ntot(\mathcal{A}^\ast) \leq \Ntot(\mathcal{A}).
\]
\end{lemma}

\begin{proof}
A CMTDs-4 winner has pairwise distinct full tuples
$(K_i,N_i,A_i,M_i)$; by determinism of $\Dec$, equal triples
$(K_i,N_i,A_i) = (K_j,N_j,A_j)$ would force $M_i = M_j$, so the triples
are already pairwise distinct and the winner is a CMTDs-3 winner on the
same output. For CMTDs-1, pairwise $\WiC_1$-distinctness makes the keys
$K_i$, hence the triples, pairwise distinct; the prechecks are
source-game rejections, and the eager length precheck is success
preserving because \TW{} decryption returns $\bot$ or a message of
length $\|C\|$, so no winning transcript is discarded. Either reduction
makes no primitive-interface query and performs at most the source-game
decryption checks, so neither cost increases.
\end{proof}

\begin{corollary}[Committing-Notion Root-Collision Bound]
\label{cor:committing-combined}
For every $s \geq 2$, every
\[
  X \in \bigl\{\mathsf{CMTDs\text{-}1},\, \mathsf{CMTDs\text{-}3},\,
              \mathsf{CMTDs\text{-}4},\, \mathsf{CMTs\text{-}1},\,
              \mathsf{CMTs\text{-}3},\, \mathsf{CMTs\text{-}4}\bigr\},
\]
and every $s$-way $X$-adversary $\mathcal{A}$ against \TW{},
\[
  \Adv^X_\TW(\mathcal{A})
  \;\leq\;
  \Lift_c(\Ntot) + \RootColl_{\tauroot}(\Nprim, s)
  \;=\;
  \Lift_c(\Ntot) + \frac{\binom{\Nprim + s}{s}}{2^{\tauroot(s-1)}}.
\]
\end{corollary}

\begin{proof}
The D-notions reduce to CMTDs-3 by
\Cref{lem:cmtds-identity-reductions} (the identity reduction for CMTDs-3
itself), whose public-permutation bound is \Cref{cor:cmtds3-pp};
monotonicity of $\Lift_c$ absorbs
$\Ntot(\mathcal{A}^\ast) \leq \Ntot(\mathcal{A})$.

The E-notions are bounded directly by the final-root
candidate-collision argument of \Cref{thm:cmtds3-ro}, with the
challenger's $s$ deterministic \emph{encryption} checks supplying the
candidate final-root transcripts in place of decryption checks. On a
winning execution all $s$ encryptions are equal, hence of equal length
and with one common root tag $T$; pairwise $\WiC_\ell$-distinctness
forces pairwise distinct triples (for $\ell = 4$ via determinism, as
above), so \Cref{cor:triple-sep} makes the $s$ final-root candidates
distinct while their $\tauroot$-bit projections all equal $T$. The win
is contained in the $s$-way candidate-collision event of
\Cref{lem:adaptive-candidate-collision}, and the flat-sponge lift
(\Cref{cor:lift-application}) transfers the bound to the
public-permutation setting at $\Nprim$; the encryption checks are
construction-internal work counted in $N$, not in $\Nprim$.
\end{proof}
